% !TEX program = lualatex
% Generated by new_lecture_note.sh
\documentclass[11pt]{article}

\usepackage{fontspec}
\usepackage{microtype}
\usepackage{mathtools,amssymb,amsfonts,amsthm,bm}
\usepackage{enumitem}
\usepackage{booktabs,array,xcolor}
\usepackage{geometry,fancyhdr,setspace}
\usepackage[hidelinks]{hyperref}
\usepackage[nameinlink,noabbrev]{cleveref}

\geometry{margin=1in,headheight=15pt}
\setstretch{1.12}
\setlength{\parindent}{0pt}
\setlength{\parskip}{0.55em plus 0.12em minus 0.08em}
\setlength{\abovedisplayskip}{0.8em plus 0.2em minus 0.1em}
\setlength{\belowdisplayskip}{0.8em plus 0.2em minus 0.1em}
\allowdisplaybreaks[3]

\setlist{
    leftmargin=*,
    topsep=0.35em,
    itemsep=0.15em,
    parsep=0pt
}

% ---------- Metadata ----------
\newcommand{\studentname}{Keshav Ramamurthy}
\newcommand{\coursename}{Math 104}
\newcommand{\courselongname}{Real Analysis}
\newcommand{\lecturenumber}{03}
\newcommand{\lecturetitle}{Lecture 03}
\newcommand{\lecturedate}{\today}

% ---------- Reusable notation ----------
\newcommand{\N}{\mathbb{N}}
\newcommand{\Z}{\mathbb{Z}}
\newcommand{\Q}{\mathbb{Q}}
\newcommand{\R}{\mathbb{R}}
\newcommand{\C}{\mathbb{C}}
\newcommand{\F}{\mathbb{F}}
\newcommand{\K}{\mathbb{K}}
\newcommand{\E}{\mathbb{E}}
\newcommand{\Prob}{\mathbb{P}}
\newcommand{\1}{\mathbf{1}}
\newcommand{\eps}{\varepsilon}
\newcommand{\dd}{\,\mathrm{d}}
\newcommand{\abs}[1]{\left\lvert#1\right\rvert}
\newcommand{\norm}[1]{\left\lVert#1\right\rVert}
\newcommand{\ip}[2]{\left\langle#1,#2\right\rangle}
\newcommand{\set}[1]{\left\{#1\right\}}
\newcommand{\given}{\,\middle\vert\,}
\newcommand{\indep}{\perp\!\!\!\perp}
\newcommand{\lto}{\longrightarrow}
\newcommand{\deriv}[2]{\frac{\mathrm{d}#1}{\mathrm{d}#2}}
\newcommand{\pderiv}[2]{\frac{\partial#1}{\partial#2}}

\DeclareMathOperator{\Aut}{Aut}
\DeclareMathOperator{\Hom}{Hom}
\DeclareMathOperator{\Ker}{ker}
\DeclareMathOperator{\im}{im}
\DeclareMathOperator{\rank}{rank}
\DeclareMathOperator{\nullity}{nullity}
\DeclareMathOperator{\Span}{span}
\DeclareMathOperator{\Var}{Var}
\DeclareMathOperator{\Cov}{Cov}

% ---------- Note environments ----------
\newtheoremstyle{notestyle}
{0.7em}{0.7em}{\itshape}{}{}{\bfseries}{.5em}{}

\theoremstyle{notestyle}
\newtheorem{theorem}{Theorem}[section]
\newtheorem{lemma}[theorem]{Lemma}
\newtheorem{proposition}[theorem]{Proposition}
\newtheorem{corollary}[theorem]{Corollary}

\theoremstyle{definition}
\newtheorem{definition}[theorem]{Definition}
\newtheorem{example}[theorem]{Example}
\newtheorem{remark}[theorem]{Remark}

\newenvironment{claim}
{\par\noindent\textbf{Claim.}\ }
{\par}

\newenvironment{idea}
{\par\noindent\textit{Idea.}\ }
{\par}

\newenvironment{proofsketch}
{\par\noindent\textbf{Proof sketch.}\ }
{\par}

% ---------- Header ----------
\pagestyle{fancy}
\fancyhf{}
\fancyhead[L]{\coursename}
\fancyhead[C]{\lecturetitle}
\fancyhead[R]{\studentname}
\fancyfoot[C]{\thepage}

\title{\vspace{-1.5em}\textbf{\coursename\ -- \lecturetitle}}
\author{\studentname}
\date{\lecturedate}

\begin{document}

\maketitle

\section{Glossary}

\begin{itemize}
    \item $A_N=\set{s_n:n\ge N}$: the \textit{tail} of the sequence, with the first
    $N-1$ terms removed.

    \item $v_N=\inf(A_N)$ and $u_N=\sup(A_N)$: the lower and upper edges of the tail.
    As the tail shrinks, $v_N$ increases and $u_N$ decreases.

    \item $\overline{\R}=\R\cup\set{\pm\infty}$: the extended reals, in which every
    monotone sequence has a limit.

    \item $s_n\to L$ means that for every $\eps>0$, there exists $N$ such that
    $\abs{s_n-L}<\eps$ for all $n\ge N$.
\end{itemize}

\section{Recall: $\limsup$ and $\liminf$}

Let $(s_n)\subset\R$ and
\[
A_N=\set{s_n:n\ge N}.
\]

Since $v_N=\inf(A_N)$ increases and $u_N=\sup(A_N)$ decreases, both converge in
$\overline{\R}$. We define
\[
\liminf s_n=\lim_{N\to\infty}\inf(A_N),
\qquad
\limsup s_n=\lim_{N\to\infty}\sup(A_N).
\]

Always,
\[
\liminf s_n\le\limsup s_n.
\]

\par
Intuitively, $\limsup s_n$ is the upper edge that the sequence approaches, while
$\liminf s_n$ is the lower edge.

\par
For example, if $s_n=(-1)^n$, every tail contains both $1$ and $-1$, so
\[
\limsup s_n=1,
\qquad
\liminf s_n=-1.
\]

\par
So when does the sequence actually converge?

\section{Convergence}

\begin{theorem}
Let $(s_n)\subset\R$. Then $(s_n)$ converges to some $L\in\R$ if and only if
\[
\liminf s_n=\limsup s_n.
\]
\end{theorem}

\par
($\implies$) Suppose $s_n\to L$. Given $\eps>0$, choose $N$ such that
\[
\abs{s_n-L}<\eps
\qquad\text{for all }n\ge N.
\]

Then every element of $A_N$ lies in $(L-\eps,L+\eps)$, so
\[
L-\eps\le v_N\le\liminf s_n
\le\limsup s_n\le u_N\le L+\eps.
\]

Since this holds for every $\eps>0$,
\[
\liminf s_n=\limsup s_n=L.
\]

\par
($\impliedby$) Suppose
\[
\liminf s_n=\limsup s_n=L\in\R.
\]

Given $\eps>0$, since $v_n\uparrow L$ and $u_n\downarrow L$, there exist
$N_1,N_2$ such that
\[
v_{N_1}>L-\eps,
\qquad
u_{N_2}<L+\eps.
\]

Let
\[
N=\max(N_1,N_2).
\]

Then for every $n\ge N$,
\[
L-\eps<v_n\le s_n\le u_n<L+\eps.
\]

Hence
\[
\abs{s_n-L}<\eps
\]
for all $n\ge N$, so
\[
s_n\to L.
\]
\section {Cauchy Sequence}
A sequence $s(n) \subset \mathbb{R}$ is Cauchy is $$\forall_{\epsilon} >0 , \exists N >
0 \text{ s.t } m, n > N \rightarrow |s_m - s_n| < \epsilon $$
\begin{lemma}[Convergence] 
\label{lem:label}
Every convergent sequence is Cauchy
\end{lemma}
\begin{proof}
    Let $(s_n) \subset \mathbb{R} $ convergent.
    Then suppose we have $\epsilon > 0$
    For $\frac{\epsilon}{2}$, $$\exists N > 0 \text{ s.t } n > N \implies |s_n - L| <
    \frac{\epsilon}{2}$$
    so its true. (just plug it in)
\end{proof}
\section{Definitions and notation}

\begin{definition}[First definition]
Write the definition and one sentence explaining why it matters.
\end{definition}

\section{Claims, theorems, and proof sketches}

\begin{claim}
State the claim in one precise sentence.
\end{claim}

\begin{proofsketch}
Record the key idea, the first useful equation, and the step that still needs checking.
\end{proofsketch}

\section{Examples and calculations}

\begin{example}[Lecture example]
Write the setup, the calculation, and the conclusion.
\end{example}

\section{Questions and follow-up}

\begin{itemize}
    \item What is still unclear?
    \item Which exercise or reading should I revisit?
    \item TODO:
\end{itemize}

\end{document}

